\section{Нормальное поле Эйлера: полный конструктивный критерий}
Пусть $R=k[Y_1,\ldots,Y_n]$, $\gamma:\A^1\hookrightarrow\A^n$ --- замкнутое вложение,
$I=I_\gamma$, $B=R[t,I/t]$. Для дифференцирования $\delta$ рассматриваем
условия
\begin{equation}\label{eq:normal}
\delta(R)\subset I,\qquad (\delta-\mathrm{id})(I)\subset I^2.
\end{equation}
Они означают исчезновение поля на кривой и единичное действие на
линейных нормальных функциях. Тогда
\[
E=\partial_t+t^{-1}\delta
\]
сохраняет $B$, поскольку $E(Y_i)=\delta Y_i/t\in B$ и
$E(F/t)=(\delta F-F)/t^2\in B$ для $F\in I$.
Локальная конечность $\delta$ означает конечномерность
$\Span_k\{\delta^jF:j\ge0\}$ для каждого $F\in R$.

\begin{theorem}\label{thm:euler}
Для заданного $\delta$, удовлетворяющего \eqref{eq:normal}, эквивалентны:
\begin{enumerate}
\item $E$ является LND на $B$;
\item $\delta$ локально конечно на $R$;
\item существуют координаты $R=k[p,q_1,\ldots,q_{n-1}]$ с
$p\gamma=u$, $I=(q_1,\ldots,q_{n-1})$ и
$\delta=\sum_jq_j\partial_{q_j}$.
\end{enumerate}
Следовательно, существование такого локально конечного поля
равносильно выпрямимости $\gamma$ над $k$.
\end{theorem}

\begin{proof}
В локализации непосредственно по индукции получаем
\begin{equation}\label{eq:falling}
E^mF=t^{-m}\prod_{j=0}^{m-1}(\delta-j)F,
\qquad
E^m(F/t)=t^{-m-1}\prod_{j=1}^{m}(\delta-j)F.
\end{equation}
Если $E$ локально нильпотентно, каждый $F$ аннулируется ненулевым
полиномом от $\delta$, что даёт локальную конечность.

Правило Лейбница и \eqref{eq:normal} дают
\[
(\delta-m)I^m\subset I^{m+1}\quad(m\ge0).
\]
Предположим теперь, что $\delta$ локально конечно.
Пусть $V$ --- конечномерное $\delta$-стабильное пространство,
$M$ --- максимальная обычная степень его элементов.
Так как $I\subset\mathfrak m_{\gamma(0)}$, имеем $V\cap I^{M+1}=0$.
Поэтому $\delta(\delta-1)\cdots(\delta-M)$ аннулирует $V$.
Его корни различны и лежат в $k$: полупростота и неотрицательные
целочисленные веса следуют автоматически, без перехода к
алгебраическому замыканию. Получаем
\[
R=\bigoplus_{d\ge0}R_d,\qquad \delta|_{R_d}=d.
\]
Для ненулевого собственного вектора точный $I$-адический порядок
равен его весу, поскольку на $I^m/I^{m+1}$ действие равно $m$.
Идеалы $I^m$ градуированы: спектральные проекторы являются
полиномами от $\delta$. Отсюда
\[
I^m=\bigoplus_{d\ge m}R_d,\quad R_0\simeq R/I\simeq k[u],\quad
R_1\simeq I/I^2,\quad R=R_0[R_1].
\]
Последнее равенство следует индукцией по весу: при $d\ge2$ имеем
$R_d\subset I^2$, и каждый однородный множитель в произведении
имеет меньший положительный вес.

Для свободности $I/I^2$ не требуется гомологическая алгебра.
Выберем $P\gamma=u$, положим $v=\gamma'$, $a=\nabla P(\gamma)$.
Тогда $av=1$, а модуль строк
$L=\{\ell\in k[u]^n:\ell v=0\}$ свободен ранга $n-1$
над PID $k[u]$. Градиент задаёт изоморфизм
\[
I/I^2\longrightarrow L,\qquad F\longmapsto\nabla F(\gamma).
\]
Явный подъём строки есть
\[
F_\ell(Y)=\sum_i\ell_i(P(Y))\bigl(Y_i-f_i(P(Y))\bigr).
\]
Инъективность доказывается Тейлором: функция с нулевым значением
и нулевым первым градиентом на кривой принадлежит $I^2$.
Выбираем инвариантный подъём $p$ параметра $u$ и базис $q_j$ модуля
$R_1$ над $R_0=k[p]$. Тогда $R=k[p,q_1,\ldots,q_{n-1}]$;
алгебраическая независимость следует из степени трансцендентности $n$.
Получена нормальная форма $\delta$.
Наконец, $B=k[t,p,q_1/t,\ldots,q_{n-1}/t]$, и $E=\partial_t$
в этих координатах. Это доказывает все импликации.
\end{proof}

\subsection{Конечная проверка и восстановление автоморфизма}
Достаточно предъявить точную стабильность
\[
V_N=\Span_k\{1,\delta^jY_i:0\le j\le N, 1\le i\le n\},\qquad
\delta^{N+1}Y_i\in V_N.
\]
Пространства произведений элементов $V_N$ стабильны по Лейбницу,
поэтому $\delta$ локально конечно на всём $R$.
Если степени элементов $V_N$ не превосходят $M$, спектральная
проекция полинома $F$ на вес $j$ вычисляется как
\[
\pr_j(F)=\prod_{\substack{0\le i\le L\\i\ne j}}
\frac{\delta-i}{j-i}F,\qquad L\ge M\deg F.
\]
Таким образом,
$p=\pr_0(P)$, $q_j=\pr_1(F_{\ell_j})$ --- искомые координаты.
Обратное отображение восстанавливается по весовым компонентам
$Y_i$ и конечным поперечным производным. Подробная формула сохранена
в \file{euler/theory.md}. Универсальной априорной границы $N$ здесь нет.

Положительный контроль восстанавливает нелинейные координаты
\[
r=z-x^2,\qquad q=y-r^2,\qquad p=x-q^2.
\]
Соответствующее пространство $V_4$ имеет размерность $9$ и точно
замыкается. Это контроль реализации, а не замена общего доказательства.

\subsection{Необходимые условия, которые ещё недостаточны}
Из нормальной формы следуют
\[
\Div\delta=n-1,\qquad (\delta Y_1,\ldots,\delta Y_n)=I.
\]
Равенство идеалов описывает всю неподвижную схему, а не только
множество точек. Постоянство дивергенции LF-поля является классическим
фактом \cite{Nowicki}; приоритет этого факта не заявляется.

Даже оба условия вместе с \eqref{eq:normal} не достаточны.
Для прямой $I=(y,z)$ поле
\[
\delta=yz\partial_x+y\partial_y+(z+y^2x)\partial_z
\]
проходит все эти тесты. Однако $\delta^{2m}x$ содержит
$xy^{3m}$ с положительным коэффициентом. Он получается повторением
пути $x\mapsto yz\mapsto y^3x$; все коэффициенты неотрицательны,
сокращений нет. Степени неограниченны, так что поле не LF.
Нельзя смешивать LF нормального поля с LND: в правильной нормальной
форме $\delta^m q_j=q_j$, а локально нильпотентным является $E$.

\section{Полная граница наивного нормального поля}
Для любого полинома параметра $P\gamma=u$ положим
\[
\delta_P=\sum_i\bigl(Y_i-f_i(P)\bigr)\partial_{Y_i}.
\]
Оно всегда удовлетворяет \eqref{eq:normal}, но имеет следующий
существенный предел применимости.

\begin{theorem}
Существует $P\gamma=u$, для которого $\delta_P$ локально конечно,
тогда и только тогда, когда существует аффинно-линейный $L$ с
$L\gamma=u$.
\end{theorem}
\begin{proof}
Если $P=L$ аффинно-линейно, то $\delta_P(P)=0$ и
$\delta_P^2Y_i=\delta_PY_i$, что даёт LF на образующих.
Обратно, необходимое равенство дивергенции имеет вид
\[
\sum_i f_i'(P)\partial_iP=1.
\]
Аффинной заменой координат приведём ненулевые положительные степени
$f_i$ к попарно различным; остальные компоненты можно сделать нулевыми.
Вид $\delta_P$ при этой замене сохраняется. Если линейного параметра
нет, все положительные степени не меньше $2$.
Пусть $d=\deg P$, а $j$ соответствует наибольшей степени $f_j$
с $\partial_jP\ne0$. Степень этого слагаемого не меньше
$(\deg f_j-1)d$, тогда как каждое предыдущее имеет степень не больше
$(\deg f_i)d-1\le(\deg f_j-1)d-1$.
Старший член не сокращается, и сумма не равна $1$.
Если все активные ненулевые компоненты отсутствуют, сумма равна нулю,
что также невозможно.
\end{proof}

Для гармонических моделей с различными индексами $a,b,c\ge2$
линейного параметра нет. Поэтому \emph{ни один} выбор $P$, какой бы
ни была его степень, не сделает это поле LF. Запрет относится также
к уже выпрямленным $H(3,4,5)$ и $H(3,5,7)$ и потому не является
препятствием выпрямимости самой кривой.

\section{Общие нормальные поля: точные пространства кандидатов}
Ищем все $\delta=a_1\partial_x+a_2\partial_y+a_3\partial_z$ с
$\deg a_i\le d$, условиями \eqref{eq:normal} и $\Div\delta=2$.
Каждый $a_i$ раскладывается по полному базису $I\cap R_{\le d}$.
При $v=\gamma'$ второе нормальное условие эквивалентно тому, что
столбцы $J_a(\gamma)-\mathrm{Id}$ кратны унимодулярному $v$.
Получается рациональная линейная система $Ac=b$, охватывающая
все поля данной степени.

\begin{center}\small
\begin{tabular}{@{}lrrrrl@{}}\toprule
Модель & $d$ & $\dim I_{\le d}$ & Неизвестных & $\operatorname{rank}A$ & Итог\\\midrule
$H(3,7,8)$ &7&63&189&189&несовместна\\
$H(3,7,8)$ &8&100&300&276&семья $24$\\
$H(4,5,7)$ &6&42&126&---&несовместна\\
$H(4,5,7)$ &7&71&213&194&семья $19$\\\bottomrule
\end{tabular}
\end{center}
Несовместность удостоверена точным рациональным вектором
$\lambda^{\mathsf T}A=0$, $\lambda^{\mathsf T}b=1$; поэтому вывод
сохраняется при любом расширении поля характеристики нуль.
Для совместных систем предъявлены рациональная точка и полный базис
направлений. Точные нулевые векторы дают верхние оценки ранга;
совпадающие модульные нижние оценки удостоверяют полноту.

Таким образом, первая степень совместности \emph{нормальных условий
вместе с глобальной дивергенцией $2$} равна $8$ и $7$.
Это не минимальные степени LF-полей или выпрямителей.
Для $H(4,5,7)$ степень $6$ невозможна уже без условия дивергенции.
Для $H(3,7,8)$ нормальные условия в степени $7$ сами по себе имеют
семейство размерности $11$; именно глобальная дивергенция исключает его.

\subsection{Два выбранных представителя строго исключены}
В семье $H(3,7,8)$ выбранный рациональный представитель сохраняет
$(x,y)$ и на факторе $k[z]$ действует как $\psi(z)\partial_z$,
$\deg\psi=6$. В семье $H(4,5,7)$ выбранный представитель сохраняет
$(y,z)$ и индуцирует $\varphi(x)\partial_x$, $\deg\varphi=7$.
Для полинома $a$ степени $e\ge2$ в характеристике нуль
\[
\deg((a\partial_v)^m v)=1+(e-1)m,\qquad
\operatorname{lc}((a\partial_v)^m v)
=\operatorname{lc}(a)^m\prod_{j=0}^{m-1}(1+(e-1)j)\ne0.
\]
Поэтому оба представителя не LF, а соответствующие $E$ не LND.
Полные коэффициенты и отдельные проверки приведены в JSON-сертификатах.
Оси не сохраняются всем семейством: соответственно $13$ из $24$ и
$11$ из $19$ выбранных базисных направлений их нарушают.
Эти числа относятся к выбранному базису; это не размерности
исключённых подсемейств. Общего запрета на все комплексные параметры нет.

\subsection{Вещественная топология и сопряжение}
По \cite[с.~2, рис.~3]{KP}, $H(3,7,8)$ имеет тип $T(2,5)$
с точностью до зеркала, а $H(4,5,7)$ --- $5_2$.
Следовательно, вещественных полиномиальных выпрямителей у этих моделей
нет. По теореме~\ref{thm:euler} \emph{любое} хорошее LF-поле для них
должно иметь невещественные коэффициенты. Это известное топологическое
ограничение поля определения, а не комплексное препятствие.

Если хорошее комплексное $\delta$ коммутировало бы с $\bar\delta$,
их среднее было бы вещественным LF-полем с теми же нормальными
условиями. Коммутирующие LF-операторы имеют конечномерные совместные
циклические пространства, поэтому среднее действительно LF.
Получилось бы запрещённое вещественное выпрямление. Аналогично,
коммутирующие $E$ и $\bar E$ дали бы вещественное LND со срезом $t$.
Следовательно, в этих случаях оба коммутатора ненулевые. При этом
\[
[E,\bar E]=t^{-2}\bigl([\delta,\bar\delta]+\delta-\bar\delta\bigr),
\]
так что два условия коммутирования нельзя отождествлять.
